\documentclass[11pt]{article}

\usepackage[T1]{fontenc}
\usepackage[margin=1in]{geometry}
\usepackage{amsmath,amssymb}
\usepackage{booktabs}
\usepackage{xcolor}
\usepackage{upquote}
\usepackage{listings}
\usepackage[hidelinks]{hyperref}
\hypersetup{
  pdftitle={nsm: an explainable, test-backed consistency engine for Standard-Model and beyond-Standard-Model gauge theories},
  pdfauthor={Evan Morritse}
}

\definecolor{codebg}{gray}{0.96}
\lstset{
  basicstyle=\ttfamily\small,
  backgroundcolor=\color{codebg},
  breaklines=true,
  frame=none,
  columns=fullflexible,
  keepspaces=true,
  xleftmargin=1em,
}

\newcommand{\nsm}{\texttt{nsm}}
\newcommand{\UY}{U(1)_Y}
\newcommand{\UX}{U(1)_X}
\newcommand{\SUthree}{SU(3)}
\newcommand{\SUtwo}{SU(2)}

\title{\nsm: an explainable, test-backed consistency engine\\
for Standard-Model and beyond-Standard-Model gauge theories}
\author{Evan Morritse}
\date{\today}

\begin{document}
\maketitle

\begin{abstract}
We describe \nsm, a symbolic-inference engine that treats the internal
consistency of a gauge theory as a constraint-satisfaction problem over an SMT
solver and elementary representation theory. From a representation skeleton
together with the gauge invariance of the Yukawa couplings and one normalization
convention, \nsm{} derives which quantum numbers are \emph{forced} (and
reports those that are non-unique, solver-undetermined, or inconsistent); classifies
which processes are allowed, forbidden, or suppressed, with a traceable reason;
evaluates hypothesized new mediators against the \emph{full} two-$U(1)$ anomaly
set, including the mixed coefficients that hand checks frequently omit; and,
given an inconsistent gauge extension, classifies the smallest \emph{tier} of
new matter that repairs it (a tiered menu from sterile singlets to
SM-vector-like exotics), with a solver-backed record (an \texttt{unsat} ladder)
that records whether smaller repairs are excluded within that menu. Every output is a structural statement,
not a numerical observable. We validate the engine by reproducing known results: the
Standard-Model hypercharges follow from anomaly cancellation and Yukawa gauge
invariance, and for gauged $B\!-\!L$ the minimal one-field sterile completion is
one right-handed neutrino per generation: rediscovered, not assumed. We emphasize what the engine does
\emph{not} do (no cross sections, no kinematics, no claim of existence) and
describe the adversarial, test-backed methodology used to check each layer.
\nsm{} is a methods/tools contribution: an explainable ``theory debugger,'' not
a calculator.
\end{abstract}

\section{Introduction}

Most computational tools for particle theory are \emph{calculators} (amplitudes,
cross sections, decay rates) or \emph{databases} (particle properties, model
files). Comparatively little tooling answers the prior, structural questions a
model builder actually asks first: \emph{Given this gauge structure, what is
forced? What is merely conventional? Is this proposed extension even
consistent, and if not, what is the smallest thing that fixes it?}

\nsm{}\footnote{\nsm{} is the package name, not an acronym; the project began
as a ``non-standard model'' pun and is now a general consistency engine for
Standard-Model and BSM gauge theories.} is an attempt to make those questions
mechanical and explainable. It
encodes gauge-theoretic consistency (anomaly cancellation, gauge invariance of
couplings, conservation laws, locality, unitarity-cut structure) as constraints
over the SMT solver Z3~\cite{z3} together with explicit representation-theory
data. The design goal is that every result be a \emph{traceable structural
statement} (``this charge is forced because dropping it leaves these anomaly
coefficients nonzero''), rather than a floating-point number whose provenance is
opaque.

\paragraph{Contributions.}
\begin{enumerate}
\item A single explanation pipeline spanning layered algebraic and graph-search engines:
\emph{detect} inconsistency, \emph{explain} the obstruction, \emph{search} a
minimal repair, \emph{re-check}, and \emph{expose} structural process
implications.
\item A two-$U(1)$ anomaly engine computing the non-abelian, gravitational,
cubic, and mixed $[Y^2X]$, $[YX^2]$ coefficients.
\item A minimal-repair search over a tiered menu (sterile singlets and
SM-vector-like, $\UX$-chiral exotic pairs) that classifies the smallest tier of
new matter an anomalous extension requires, with a solver-backed record
reporting minimality within that menu: a consistency checker that doubles as a
modest designer.
\item Validation by reproduction of known results and by adversarial tests in
which each layer is independently re-derived.
\end{enumerate}

\paragraph{Non-contributions.}
\nsm{} computes no physical observables and makes no claim of discovery. It does
not derive masses, mixing angles, coupling values, or the number of generations
(these remain empirical inputs). Its results are structural and
quantum-number-level; see Section~\ref{sec:limits}. Throughout we use the
convention $Q = T_3 + Y/2$, so the quark doublet carries $Y = 1/3$.

\section{Design}

\nsm{} is organized as layers over one explanation pipeline
(Table~\ref{tab:layers}). Lower layers expose primitives (anomaly coefficients,
a tree-level reachability search) that upper layers reuse, so that, e.g., a
candidate gauge boson's consistency check is \emph{literally} the kernel's
anomaly-cancellation routine, and the process-mediation test is a special case
of factorization-channel enumeration.

\begin{center}
\centering
\small
\begin{tabular}{@{}lp{0.62\linewidth}@{}}
\toprule
Layer & Establishes \\
\midrule
kernel & hypercharges/charges \emph{forced} by anomaly cancellation + Yukawa
gauge invariance; reports unique / non-unique / solver-undetermined / inconsistent \\
dynamics & layered process verdict from conservation laws + a tree-level vertex
search \\
extensions & field type of a candidate mediator (implemented coupling templates);
consistency (gauge candidates reuse the anomaly engine); process exposure \\
amplitudes & leading soft factors and their Ward identities \\
factorization & internal lines a process factorizes through (locality) \\
unitarity & intermediate states a process can pass through (optical-theorem
cuts), ordered by mass threshold \\
repair & smallest addition that makes an inconsistent gauge extension consistent \\
\bottomrule
\end{tabular}
\par\medskip
\refstepcounter{table}\label{tab:layers}
Table~\thetable: The layers of \nsm. Each is independently re-derived and tested.
\end{center}

\section{Related work}

\nsm{} is closest in spirit to model-building infrastructure, but it is aimed
at a different phase of the workflow. FeynRules~\cite{feynrules2} and
SARAH~\cite{sarah4} take a specified model (fields, symmetries, parameters,
and often a Lagrangian) and derive model data such as vertices, mass matrices,
renormalization-group equations, and interfaces to simulation tools. PyR@TE
focuses on automatic renormalization-group equations for general gauge
theories~\cite{pyrate}. SModelS decomposes BSM collider signatures into
simplified-model topologies and confronts them with LHC limits~\cite{smodels}.
These tools are deeper calculators and phenomenology interfaces than \nsm{} is.

The anomaly-equation literature addresses closely related mathematics, including
general solutions of chiral $U(1)$ anomaly equations~\cite{u1solution}. \nsm{}
does not compete with those classifications. Its contribution is instead an
explainable inverse workflow: diagnose why a proposed extension is inconsistent,
search a declared repair menu, record why smaller repairs fail, re-check the
result, and connect the repaired mediator to structural process exposure. This
is a deliberately narrower, more inspectable layer that can sit before a full
model is handed to a calculator such as FeynRules or SARAH.

\section{Forced quantum numbers}
\label{sec:kernel}

Electric charge is not a primitive input: it follows from weak isospin and
hypercharge through $Q = T_3 + Y/2$. The kernel leaves all hypercharges
\emph{unknown} and imposes consistency as constraints:
\begin{itemize}
\item anomaly cancellation: the non-abelian, mixed, gravitational, and cubic
coefficients must vanish, summed over Weyl fermions with a chirality sign
$s=\pm 1$;
\item gauge invariance of each Yukawa coupling (so fermions can acquire mass);
\item one normalization ($Y_H = 1$).
\end{itemize}
Z3 then returns the hypercharges. For one Standard-Model generation the solution
is unique:
\begin{equation}
Y_{Q_L}=\tfrac13,\quad Y_{u_R}=\tfrac43,\quad Y_{d_R}=-\tfrac23,\quad
Y_{L_L}=-1,\quad Y_{e_R}=-2,
\end{equation}
whence $Q_u=+\tfrac23$, $Q_d=-\tfrac13$, $Q_\nu=0$, $Q_e=-1$.

The value of the constraint formulation is that it equally reports what is
\emph{not} forced. Dropping the Yukawa couplings leaves an exact two-fold
$u_R\!\leftrightarrow\!d_R$ hypercharge ambiguity; adding a free right-handed
neutrino opens a one-parameter flat direction (gauged $B\!-\!L$); a lone colored
Weyl fermion is rejected because $[\SUthree^3]\neq 0$; an odd number of weak
doublets is rejected by the Witten $\SUtwo$ global anomaly~\cite{witten}. The
engine distinguishes \emph{forced}, \emph{free}, \emph{solver-undetermined}, and
\emph{inconsistent} as first-class outcomes.

\section{Processes and structural amplitude constraints}

Above the kernel, a process $i\to f$ receives a layered verdict from (a) exact
gauge symmetries (charge, color), (b) the accidental global symmetries of the
renormalizable theory (baryon and lepton number), (c) approximate symmetries
(lepton flavor), (d) kinematics, and (e) a tree-level vertex-reachability search.
The search makes ``$\mu^-\!\to e^-\bar\nu_e\nu_\mu$ proceeds at tree level via
$W$ exchange'' a \emph{computed} statement rather than a conservation tautology,
and correctly finds no tree diagram for $\mu^-\!\to e^-\gamma$.

Three further layers encode \emph{amplitude-inspired} structural constraints:
the quantum-number content of $S$-matrix consistency conditions, without
introducing kinematics or evaluating amplitudes. The soft-theorem layer records
the leading soft factors~\cite{weinbergsoft} and their Ward identities (a
spin-1 emission is consistent iff the signed charge flow vanishes; a spin-2
emission iff the coupling is universal, reducing the Ward identity to momentum
conservation). The factorization layer enumerates the internal lines through
which a process factorizes (locality). The unitarity layer enumerates the
multi-particle intermediate states a process can pass through (optical-theorem
cuts), ordered by mass threshold. Each captures the structural content of an
$S$-matrix consistency condition (charge conservation under soft emission,
locality, the optical theorem) at the level of vertices and quantum numbers;
none evaluates an amplitude, a residue, or a phase-space integral.

\section{Candidate extensions and the full anomaly set}

Within the flat-space, Lorentz-invariant, long-range mediator setting modeled
here, a hypothesized new mediator is classified by the rank of the conserved
current it couples to: a spin-$s$ field couples to a rank-$s$ current (a
classification heuristic within the modeled class, not a theorem about all
mediators), while the implemented coupling templates in this prototype are a
scalar source (rank $0$), a conserved gauge current (rank $1$), and the
stress-energy tensor (rank $2$)
\cite{colemanmandula,weinbergwitten}. A massless mediator coupling universally
to energy--momentum is therefore classified as spin~2. We treat massless
higher-spin long-range couplings as outside the consistency class modeled here,
rather than as targets for this structural search.

For a newly \emph{gauged} $\UX$ alongside hypercharge, the full anomaly set
includes pure-$X$ and \emph{mixed} coefficients. Writing $d_3,d_2$ for the color
and weak dimensions, $T_c,T_w$ for the Dynkin indices, $A_c$ for the $\SUthree$
cubic coefficient, and $s=\pm1$ for chirality, the engine evaluates
\begin{align}
[\SUthree^2 X] &= \textstyle\sum s\,d_2 T_c\,X, &
[\SUtwo^2 X] &= \textstyle\sum s\,d_3 T_w\,X, &
[\text{grav}^2 X] &= \textstyle\sum s\,d_3 d_2\,X, \nonumber\\
[X^3] &= \textstyle\sum s\,d_3 d_2\,X^3, &
[Y^2 X] &= \textstyle\sum s\,d_3 d_2\,Y^2 X, &
[Y X^2] &= \textstyle\sum s\,d_3 d_2\,Y X^2,
\label{eq:mixed}
\end{align}
together with the corresponding $Y$-sector and $[\SUthree^3]$ coefficients. The
mixed terms~\eqref{eq:mixed} are load-bearing: a charge assignment can cancel
every pure-$X$ anomaly yet fail $[Y^2 X]$, so a single-$U(1)$ check can certify
a still-anomalous model as consistent.

\section{Minimal repair}

Given a charge assignment that fails the conditions above, \nsm{} searches for
the smallest set of added fermions whose charges cancel the residual anomalies.
The anomaly coefficients are linear, quadratic, and cubic in the unknown added
charges, so the search is naturally posed for the SMT solver.

The menu is \emph{tiered}. The simplest entry is the \emph{sterile right-handed
singlet} (color and weak singlet with $Y=0$ and unknown $X$): neutral under
everything except $\UX$, such a field contributes only to $[\text{grav}^2 X]$
and $[X^3]$ (each $\propto -X$ and $-X^3$ respectively), so cancelling a residual
$(g,c)=([\text{grav}^2 X],[X^3])$ requires
\begin{equation}
\sum_i X_i = g, \qquad \sum_i X_i^3 = c,
\label{eq:repair}
\end{equation}
solved for gauged $B\!-\!L$ by a single $\nu_R$. Beyond the singlet, the menu
adds \emph{SM-vector-like, $\UX$-chiral} exotic pairs $E,L,D,Q$ in definite SM
representations with left/right charges $X_L\neq X_R$. Being vector-like under
the Standard Model, each pair cancels its own SM anomalies; its chiral $\UX$
split feeds the $X$-sector, with contributions $\propto(X_L-X_R)$ and
$(X_L^3-X_R^3)$. A hypercharged pair ($E$ or $L$) thereby reaches the mixed
$[Y^2X]$, $[YX^2]$, and a colored pair ($D$ or $Q$) reaches $[\SUthree^2 X]$:
precisely the anomalies a sterile singlet cannot touch. The chirality condition
$X_L\neq X_R$ is essential; an $X$-vector-like pair would be inert.

The engine returns the minimum \emph{tier} that repairs the model (sterile
singlet, then colorless exotic, then colored exotic), at minimal field count and
then representation complexity. If a nonzero anomaly lies outside the reach of
every menu item, the model is reported \emph{structurally blocked} with that
coefficient named; if every nonzero anomaly is reachable but no repair fits the
field budget, it is reported \emph{budget-limited} instead. The tool thus never
claims a fix it cannot deliver, nor mistakes a budget limit for an obstruction.

The search emits a \emph{solver-backed minimality record} rather than merely
asserting minimality: for every field count below the one returned, the solver's
verdict is recorded. An \texttt{unsat} verdict is a decision over the reals (no charges of
that size satisfy~\eqref{eq:repair}, rational or otherwise), so a ladder of
\texttt{unsat} rungs certifies that no smaller repair exists within the menu. A
size the solver cannot decide (\texttt{unknown}), or one for which the solver
returned a nonrational (algebraic) model (which does not prove that no
rational solution exists), is recorded honestly and minimality is \emph{not} claimed.
This record's soundness rests only on that of the solver's
nonlinear-real-arithmetic decision procedure (no false \texttt{unsat}), so
``minimal within the menu'' becomes a recorded result rather than an assertion.
(We record the solver's verdicts, but not tracked unsat cores, SMT-LIB replay
artifacts, or proof logs. An unsat core would identify a conflicting subset of
assumptions; a proof log, where supported, would provide the stronger proof
artifact.) For gauged $B\!-\!L$ it is immediate: the only smaller content is the
unrepaired, inconsistent model, so a single $\nu_R$ is minimal.

The reported repair is minimal \emph{within the menu and the field-count cost};
it is not unique, not a naturalness statement, and not a claim that the field
exists.

\section{\texorpdfstring{Case study: gauging $B\!-\!L$}{Case study: gauging B-L}}
\label{sec:casestudy}

Assign each fermion its $B\!-\!L$ charge ($+1/3$ for quarks, $-1$ for leptons)
and gauge it. The full anomaly set (Table~\ref{tab:bml}) shows the model is
inconsistent, and \emph{why}: only $[\text{grav}^2 X]$ and $[X^3]$ fail, both
equal to $-1$; every mixed and non-abelian coefficient already cancels.

\begin{table}[t]
\centering
\small
\begin{tabular}{@{}lrlr@{}}
\toprule
coefficient & value & coefficient & value\\
\midrule
$[\SUthree^3]$ & $0$ & $[\text{grav}^2 X]$ & $-1$ \ \textbf{(fails)}\\
$[\SUthree^2 Y]$, $[\SUthree^2 X]$ & $0$ & $[Y^3]$ & $0$\\
$[\SUtwo^2 Y]$, $[\SUtwo^2 X]$ & $0$ & $[Y^2 X]$, $[Y X^2]$ & $0$\\
$[\text{grav}^2 Y]$ & $0$ & $[X^3]$ & $-1$ \ \textbf{(fails)}\\
\bottomrule
\end{tabular}
\caption{The eleven anomaly coefficients for gauged $B\!-\!L$ with
Standard-Model fermions only, as computed by \nsm. Only the gravitational and
cubic $X$-anomalies fail.}
\label{tab:bml}
\end{table}

By~\eqref{eq:repair} with $g=c=-1$, the unique single-field repair is $X=-1$:
one sterile right-handed singlet of $B\!-\!L$ charge $-1$, i.e.\ a right-handed
neutrino. Re-running the full anomaly set on the repaired content returns all
zeros. The engine has rediscovered the standard motivation for $\nu_R$ in gauged
$B\!-\!L$~\cite{bminusL}, without being told the answer. With $\nu_R$ in place
the $B\!-\!L$ $Z'$ is consistent, and the exposure layer reports it as resonantly
produced in $e^+e^-\!\to Z'$, as an $s$-channel resonance/interference
contribution to $e^+e^-\!\to\mu^+\mu^-$, and, being flavor-diagonal, as
\emph{not} a source of
$\mu\to e Z'$. This single run threads gauged $B\!-\!L$, the full anomaly check,
the minimal sterile repair, the $\nu_R$ motivation, and $Z'$ phenomenology into
one explained chain.

Gauged $B\!-\!L$ is only the sterile corner of the repair map. The same engine
classifies the \emph{minimum tier} of new matter required for any $\UX$ charge
assignment (Table~\ref{tab:tiers}), and a second worked case makes the tiering
concrete. Assign $X=1$ to $e_R$ alone; the engine returns the anomaly vector
\[
\bigl([\text{grav}^2X],\,[Y^2X],\,[YX^2],\,[X^3]\bigr) = (-1,\,-4,\,+2,\,-1),
\qquad [\SUthree^2X]=[\SUtwo^2X]=0 .
\]

A sterile singlet has $Y=0$, so it feeds only $[\text{grav}^2X]$ and $[X^3]$ and
\emph{cannot} touch the mixed $[Y^2X]$ or $[YX^2]$; the engine reports the
sterile menu as structurally blocked on exactly those two coefficients. The
search is therefore forced up one tier, and the colorless menu repairs the model
with one hypercharged $E_L/E_R$ pair ($Y=-2$): two Weyl fields carrying
$X_L=1$, $X_R=0$ in the signed left/right convention of Eq.~\eqref{eq:mixed},
vector-like under the Standard Model and chiral under $\UX$. With that pair
included, all eleven coefficients re-verify to zero.

This example isolates anomaly repair. Invariance of the ordinary electron
Yukawa interaction under the added $\UX$ is a separate model-building constraint
unless the scalar charges or scalar sector are included in the skeleton.

The minimality record is a genuine ladder: the one-field ($1N$) and two-field
($2N$) sterile attempts are both \texttt{unsat} (no lone singlet, and no pair of
singlets, can generate the $[Y^2X]$ it must cancel), so the two-field colorless
pair is certified minimal within the declared menu and cost function.
A charge on $Q_L$ alone instead breaks the colored $[\SUthree^2X]$, which only
$D$ or $Q$ can reach, pushing the same machinery one tier further.

\begin{table}[t]
\centering
\small
\begin{tabular}{@{}lll@{}}
\toprule
charge assignment & minimum repair tier & example repair\\
\midrule
$B\!-\!L$ & sterile-repairable & $\nu_R$ with $X=-1$\\
$X$ on $e_R$ only & colorless-exotic-repairable & an $E_L/E_R$ chiral split\\
$X$ on $Q_L$ only & colored-exotic-required & a $Q_L/Q_R$ chiral split\\
\bottomrule
\end{tabular}
\caption{The minimum repair tier for three $\UX$ charge assignments, as
classified by \nsm{}. $B\!-\!L$ is the familiar sterile case; the others demand
hypercharged or colored exotics that the sterile menu cannot supply.}
\label{tab:tiers}
\end{table}

This does not discover particles; it systematically classifies the smallest
completion within a declared representation menu and cost function. That turns
the layer from a checker that validates one known result into an explorer that
maps a space of anomalous extensions into minimal-completion classes, reporting
what \emph{kind} of new matter a broken gauge extension minimally demands.

\paragraph{A repair census.}
Because the classifier is mechanical, we can run it over a family of extensions
rather than a single example. We enumerate every integer $\UX$ charge assignment
to one generation with $|q|\le 2$ and reduce rescaling/sign duplicates to
primitive representatives (the tier is invariant under $X\to\lambda X$; it is
also invariant under hypercharge mixing $X\to X+\alpha Y$, which we do not
quotient but use as an internal consistency check), then classify each by its
cheapest repair tier. The resulting census (Table~\ref{tab:census}) is a
reproducible, basis-dependent slice: of the $1442$ assignments distinct under
sign and overall integer rescaling (one of which is the trivial zero generator,
the lone ``already-consistent'' cell and not itself a gauged extension), $0.1\%$
are sterile-repairable, $9.6\%$ need only colorless exotics, $31.5\%$ require
colored exotics, and the remaining $58.7\%$ are not repaired within a six-field
budget.

\begin{table}[t]
\centering
\small
\begin{tabular}{@{}lrr@{}}
\toprule
minimum repair tier & count & share\\
\midrule
already-consistent & $1$ & $0.1\%$\\
sterile-repairable & $1$ & $0.1\%$\\
colorless-exotic-repairable & $139$ & $9.6\%$\\
colored-exotic-required & $454$ & $31.5\%$\\
budget-limited ($>6$ fields) & $847$ & $58.7\%$\\
\bottomrule
\end{tabular}
\caption{Repair census of all $1442$ sign- and rescaling-distinct
(primitive, sign-fixed) integer $\UX$ assignments to one generation with $|q|\le2$ in the
fixed $(Y,X)$ basis, each classified by the cheapest repair tier within a budget
of six added fields (menu N/E/L/D/Q, solver timeout $8\,$s). The raw grid has
$3125$ cells ($2.17\times$ deduplication); $0$ cells are structurally blocked and
$81$ ($5.6\%$) carry an uncertified solver rung.}
\label{tab:census}
\end{table}

Three caveats are built into the table. First, the budget-limited majority is a
\emph{field-budget artifact}, not an unrepairability theorem: with the full menu
every failing anomaly is individually reachable (a necessary, not sufficient,
condition for a joint repair), no cell is structurally blocked, and the
budget-limited fraction falls monotonically as the budget grows (it has not
plateaued at six fields). We therefore report budget-limited as ``not found
within the budget,'' not as a proof that a repair exists at larger size. Second, the percentages count a bounded integer box
in one fixed $(Y,X)$ basis; they are not a measure over all $\UX$ extensions, and
a different choice of generator would re-slice the population. Third, $5.6\%$ of
cells carry an uncertified solver rung (an \texttt{unknown} or
\texttt{sat-nonrational-model} verdict in the minimality record), which we report rather
than suppress. With those stated, the census is the method claim made concrete:
not one rediscovery, but a reproducible map of which kind of new matter a space
of anomalous gauge extensions minimally demands.

\section{Implementation}
\label{sec:impl}

\nsm{} is implemented in Python over the Z3 SMT solver as focused modules
sharing one engine; the entire narrated pipeline runs via
\texttt{python -m nsm}. Consistency conditions are compiled to solver
assertions, and the routine that \emph{checks} a concrete spectrum is the same
one that \emph{solves} for an unknown one: the anomaly-coefficient arithmetic is
polymorphic in its charge arguments, returning exact rationals when fed
rationals (a check) and symbolic expressions when fed solver variables (a
constraint). Listing~\ref{lst:api} shows a representative query.

\begin{lstlisting}[language=Python,caption={Public API: forced hypercharges, and a minimal repair.},label=lst:api]
from fractions import Fraction as F
from nsm import derive_hypercharges, minimal_repair
from nsm.sm import standard_model_skeleton

r = derive_hypercharges(standard_model_skeleton())
r.hypercharges["Q_L"], r.is_unique     # -> (Fraction(1, 3), True)

# gauged B-L (charges as exact Fractions); repair searches for a fix
B_minus_L = (
    ("Q_L", F(1, 3)), ("u_R", F(1, 3)), ("d_R", F(1, 3)),
    ("L_L", F(-1)), ("e_R", F(-1)),
)
minimal_repair(B_minus_L).added        # -> [('nu_R', Fraction(-1, 1))]
\end{lstlisting}

The repair-census layer uses the same primitive over an integer charge grid: it
reduces each vector to a primitive sign-fixed representative, keeps the raw and
deduplicated denominators, classifies the representative by its minimum repair
tier, and exports CSV, Markdown, or \LaTeX{} summaries. Hypercharge-basis shifts
are treated as an equivalence check rather than as a counting quotient, since
they shear the integer lattice.

The kernel poses the hypercharge problem as a conjunction of assertions (each
anomaly coefficient equal to zero, gauge invariance of each Yukawa coupling, and
one normalization) and reads the unique assignment from the model. Whether that
assignment \emph{is} unique is decided by a second solve: assert the negation of
the found solution and check satisfiability. \texttt{unsat} certifies
uniqueness; \texttt{sat} exhibits a distinct second solution (such as the
$u_R\!\leftrightarrow\!d_R$ ambiguity of Section~\ref{sec:kernel}); and
\texttt{unknown} (the solver exceeding its budget on the nonlinear real
arithmetic) is reported as a genuine third outcome rather than silently
collapsed to false certainty. The same reachability and satisfiability idioms
serve the dynamics, factorization, and unitarity layers. The minimal-repair
search (Listing~\ref{lst:repair}) encodes anomaly cancellation for each
candidate menu combination, increasing the field count until the smallest
solution is found and skipping nonrational solver models rather than failing on
them.

\begin{lstlisting}[language=Python,caption={Tiered repair: search the declared menu, smallest cost first.},label=lst:repair]
for fields, complexity, combo in get_combinations(max_added, menu):
    if combo_cannot_reach_a_nonzero_anomaly(combo):
        record(combo, "unsat")          # sound prune, no solver call
        continue
    status, charges = solve_anomaly_equations(combo)
    record(combo, status)
    if status == "sat":
        repair = make_added_fields(combo, charges)
        assert verify_all_anomalies_vanish(repair)
        return repair                  # first sat is minimal in this menu/cost
\end{lstlisting}

\section{Validation}

Because the engine's value is its correctness, each layer was checked two ways.
First, a test suite (171 tests, distributed with the code) pins the physics
against hand-computed values, including \emph{nonzero} anomaly coefficients, so
that a wrong Dynkin or multiplicity factor which still cancels for the symmetric
Standard Model cannot hide. Second, as internal quality assurance, each layer's
results were independently re-derived and differential-tested against the
implementation over randomized inputs; this process found and fixed real
defects (for example, a non-terminating uniqueness search and a crash on
nonrational solver models) and indicated that the full two-$U(1)$ check agrees
with a single-$U(1)$ check wherever the latter applies while additionally
catching mixed-anomaly failures, with no disagreement observed over
$2\times10^{4}$ random charge assignments. This internal cross-checking is
evidence of correctness, not a substitute for independent external reproduction;
the test suite is the reproducible artifact.

\paragraph{Guaranteed versus observed.}
The engine rests on two trust anchors: anomaly coefficients for concrete
spectra are computed in exact rational arithmetic, and the nontrivial solver
inference is delegated to Z3's nonlinear-real-arithmetic decision procedure. A
reported zero or nonzero coefficient is therefore exact, and an \texttt{unsat}
verdict rules out a repair of that size over the reals. A ladder of
all-\texttt{unsat} rungs is consequently a solver-backed guarantee of
minimality within the declared menu and cost function. Everything else is
reported as evidence, not proof: an \texttt{unknown} or \texttt{sat-nonrational-model}
rung does not certify minimality, and cross-layer differential agreement is a
test result rather than a theorem. For the anomaly and repair layers the tool
either returns a solver-backed guarantee or labels the result uncertified; the
structural layers (process, soft-factor, factorization, unitarity) are heuristic
within their modeled class.

\section{Limitations}
\label{sec:limits}

\nsm{} is a structural / quantum-number engine. It computes no cross sections,
decay rates, or branching ratios; it has no kinematics (no momenta, spinors, or
phase-space integrals; mass enters only as a crude threshold); it evaluates no
amplitude values or pole residues; it performs no detector simulation and makes
no statement of statistical significance or discovery. It does not derive
masses, mixing angles, coupling constants, or the number of generations, which
remain empirical inputs; an apparent ``prediction'' of one of these would be a
bug. The vertex set is CKM-diagonal, composite hadrons are not expanded into
constituents, and candidate gauge couplings are taken family-universal and
vector-like. Several layer outputs are deliberately narrower than their names
suggest: ``color'' conservation means \emph{color triality} (a necessary, not
sufficient, color-singlet test); baryon and lepton number refer to the modeled
renormalizable vertex set; ``kinematics'' means decay thresholds and an
$n\!\to\!1$ phase-space check, not general momentum-space kinematics; and the
unitarity layer enumerates \emph{candidate} cut topologies with nominal mass
thresholds, not cuts known to be open at a given $\sqrt{s}$. The one-line
summary: \nsm{} establishes \emph{structural and
quantum-number consistency}; it does not compute, measure, or predict any
physical observable.

\section{Conclusion}

\nsm{} demonstrates that a useful fraction of gauge-theory reasoning (what is
forced, what is free, what is forbidden, and how to repair an inconsistent
extension) can be made automatic, explainable, and mechanically checked, with
a clean boundary between ``structurally consistent'' and ``actually computed.''
Natural extensions include emitting tracked unsat cores, SMT-LIB replay
artifacts, or solver proof logs for the minimality records, and, as a distinct
phase that would cross the structural boundary, a spinor-helicity layer enabling
genuine residue factorization and amplitude values.

\paragraph{Availability.} The implementation, test suite, and a worked
case study are in the accompanying repository.

\begin{thebibliography}{9}
\small
\bibitem{z3} L.~de Moura and N.~Bj{\o}rner, \emph{Z3: An Efficient SMT Solver},
in TACAS 2008, LNCS \textbf{4963} (Springer), pp.~337--340 (2008).
\bibitem{feynrules2} A.~Alloul, N.~D.~Christensen, C.~Degrande, C.~Duhr, and
B.~Fuks, \emph{FeynRules 2.0--A complete toolbox for tree-level
phenomenology}, Comput.\ Phys.\ Commun.\ \textbf{185}, 2250--2300 (2014).
\bibitem{sarah4} F.~Staub, \emph{SARAH 4: A tool for (not only SUSY) model
builders}, Comput.\ Phys.\ Commun.\ \textbf{185}, 1773--1790 (2014).
\bibitem{pyrate} F.~Lyonnet, I.~Schienbein, F.~Staub, and A.~Wingerter,
\emph{PyR@TE: Renormalization Group Equations for General Gauge Theories},
Comput.\ Phys.\ Commun.\ \textbf{185}, 1130--1152 (2014).
\bibitem{smodels} S.~Kraml, S.~Kulkarni, U.~Laa, A.~Lessa, W.~Magerl,
D.~Proschofsky-Spindler, and W.~Waltenberger, \emph{SModelS: a tool for
interpreting simplified-model results from the LHC and its application to
supersymmetry}, Eur.\ Phys.\ J.\ C \textbf{74}, 2868 (2014).
\bibitem{u1solution} D.~B.~Costa, B.~A.~Dobrescu, and P.~J.~Fox,
\emph{General Solution to the $U(1)$ Anomaly Equations}, Phys.\ Rev.\ Lett.\
\textbf{123}, 151601 (2019).
\bibitem{witten} E.~Witten, \emph{An $SU(2)$ Anomaly}, Phys.\ Lett.\ B
\textbf{117}, 324 (1982).
\bibitem{weinbergsoft} S.~Weinberg, \emph{Infrared Photons and Gravitons},
Phys.\ Rev.\ \textbf{140}, B516 (1965).
\bibitem{colemanmandula} S.~Coleman and J.~Mandula, \emph{All Possible
Symmetries of the $S$ Matrix}, Phys.\ Rev.\ \textbf{159}, 1251 (1967).
\bibitem{weinbergwitten} S.~Weinberg and E.~Witten, \emph{Limits on Massless
Particles}, Phys.\ Lett.\ B \textbf{96}, 59 (1980).
\bibitem{bminusL} R.~N.~Mohapatra and R.~E.~Marshak, \emph{Local $B\!-\!L$
Symmetry of Electroweak Interactions, Majorana Neutrinos and Neutron
Oscillations}, Phys.\ Rev.\ Lett.\ \textbf{44}, 1316 (1980).
\end{thebibliography}

\end{document}
